એક \emph{ફલન} $f\colon A\to B$ તેના \emph{પ્રદેશ} $A$ના દરેક ઘટકને \emph{સહપ્રદેશ} $B$ના એક જ ઘટક સાથે જોડે છે. $x\in A$ હોય ત્યારે $f$ તેને જે $y$ સાથે જોડે તેને \emph{દલીલ} $x$ માટે $f$નું \emph{મૂલ્ય} $f(x)$ કહીએ છીએ. જો $A$ જોડોનો ગણ હોય તો $f$ બે દલીલો લે છે એમ માની શકીએ. $f$નો \emph{વિસ્તાર} $\ran{f}$ એ $B$નો એવો ઉપગણ છે જેમાં $f$નાં બધાં મૂલ્યો છે.

જો $\ran{f}=B$ હોય તો $f$ \emph{વ્યાપ્ત} છે. $f$ એક $x$ને એક જ $f(x)$ સાથે જોડે છે. જો બે જુદી દલીલોનાં મૂલ્યો હંમેશાં જુદાં હોય તો $f$ \emph{એક-એક} છે. એક-એક અને વ્યાપ્ત ફલનને \emph{એક-એક વ્યાપ્ત} કહે છે.

એક-એક વ્યાપ્ત ફલનનું એકમાત્ર \emph{પ્રતિફલન} $f^{-1}$ હોય છે. ફલનોને સાંકળી પણ શકાય છે. $(g\circ f)$ એ $f$ પછી $g$નું \emph{સંયોજન} છે. એક-એક ફલનોનું સંયોજન એક-એક હોય છે અને વ્યાપ્ત ફલનોનું સંયોજન વ્યાપ્ત હોય છે. $(f^{-1}\circ f)$ તદેવ ફલન છે.

$f$નું દરેક $x\in A$ પર મૂલ્ય હોવાની શરત છોડીએ તો \emph{આંશિક ફલન} મળે છે. જો $f\colon A\pto B$ આંશિક હોય અને દલીલ $x$ માટે $f$નું મૂલ્ય હોય તો $f(x)$ \emph{વ્યાખ્યાયિત} છે, $f(x)\fdefined$. અન્યથા તે \emph{અવ્યાખ્યાયિત} છે, $f(x)\fundefined$. દરેક (આંશિક) ફલન $f$નો \emph{આલેખ} $R_f$ એ એવો સંબંધ છે જે ત્યારે અને માત્ર ત્યારે લાગુ પડે છે જ્યારે $f(x)=y$.
